\section{Introduction}

The stationary equilibrium of \citet{Aiyagari1994} is a number: the interest rate at which the
capital that a continuum of households accumulates against uninsurable earnings risk equals the
capital a competitive firm demands. Aiyagari computed it by bisection on the rate and reported it,
for three degrees of relative risk aversion $\mu\in\{1,3,5\}$ and eight earnings processes, in his
Table II. He did not ask whether the number was unique, and for three decades the question was
treated as settled by the picture: capital supply is drawn rising in the rate, demand falling, and
the curves cross once.

The picture is not a theorem. \citet{Acikgoz2018} proved existence and gave conditions for
uniqueness that do not cover Aiyagari's parameters; \citet{Light2020} proved uniqueness for
relative risk aversion at most one, again under conditions on the discount factor and the earnings
process that his own statement does not verify at $\beta=0.96$; \citet{Light2022} put the argument
in a general mean-field framework; \citet{Kirkby2019} computed the Aiyagari equilibrium with a discretisation proved to converge to
the true solution, found numerically that it is almost certainly unique, and recorded that
uniqueness remained open theoretically; \citet{YoungWalsh2025} showed by computation that at high
risk aversion the supply curve can bend backward and the equilibrium can fail to be unique. The survey
of \citet{TodaWalsh2024} lists uniqueness in the Aiyagari model with its actual calibration as
open.

This paper closes it for $\mu=1$ and locates the obstruction for $\mu\in\{3,5\}$, and it does so in
a form that is new to economics: every result is checked by a proof assistant. The whole
development, from the definition of the household problem to the theorem that two equilibrium rates
of Aiyagari's log economy coincide, is written in Lean 4 \citep{Lean4} against the Mathlib library
\citep{Mathlib}, and the Lean kernel has verified that each proof follows from its hypotheses and
the standard axioms of the library. There are no unproved steps, no appeals to results that are
``well known'', and no numerical approximations. The formal statement of each result reported below
is named in the text, and Appendix~\ref{app:index} lists them with the files that contain them.

\paragraph{The main result.}
Fix Aiyagari's preferences and technology at $\mu=1$: log utility, $\beta=0.96$, Cobb--Douglas
production with capital share $\alpha=0.36$ and depreciation $\delta=0.08$, no borrowing. Let the
labour endowment follow any finite-state Markov chain whose transitions are monotone (a higher
state's row first-order stochastically dominates a lower state's) and whose lowest state is reached
from every state in one step. Then at most one interest rate in $(-\delta,\lambda)$, with
$\lambda=1/\beta-1$, is a stationary equilibrium rate (Theorem~\ref{thm:main}). Aiyagari's own
earnings processes, Tauchen discretisations of an AR(1) with positive persistence on any number of
grid points, satisfy the two conditions by construction (Corollary~\ref{cor:tauchen}). Outside
$(-\delta,\lambda)$ there is nothing: above $\lambda$ the household is patient and accumulates
without bound, so no stationary distribution has finite mean capital below demand
(Theorem~\ref{thm:above}); below $-4\%$ the only capital a log household holds is bounded by
$\beta\bar y/(1-\beta(1+r))$, which falls under demand as demand explodes at $-\delta$
(Theorem~\ref{thm:below}). Existence on any interval is reduced to a single inequality on capital
supply at its upper end (Theorem~\ref{thm:exist}), and the bisection Aiyagari ran is proved to
bracket and converge to the unique rate (Theorem~\ref{thm:bisect}).

\paragraph{Why the published arguments do not apply, and what does.}
The economics of the proof is the chain of \citet{Light2020}: individual saving rises with the
rate (his Theorem 1); therefore aggregate capital, the mean of the stationary distribution, rises
with the rate (his Theorem 2); therefore supply meets the strictly decreasing demand at most once
(his Theorem 3). Each link needed something new at $\beta=0.96$.

Theorem 1 is where relative risk aversion enters, through the monotonicity of $c\,u'(c)$ that the
rescaling argument needs, and for log utility that is available. But the argument is an induction
along value-function iterates, and it needs the borrowing constraint to be slack at the states it
compares and the asset cap, an artefact of compactness, to be slack everywhere. The published
sufficient conditions for both involve Lipschitz constants of the value function that scale like
$1/(1-\beta R)$, which at Aiyagari's rates is several hundred; at his parameters the tests fail by
orders of magnitude. What replaces them are two consequences of the Euler inequality itself: a
household at the constraint consumes at least the lowest income, which turns the corner condition
into the primitive inequality $\beta R\,u'(y_{\min})<u'(\text{cash on hand})$
(Section~\ref{sec:corner}); and a minimal marginal propensity to consume, $\kappa=1-\beta$ for log
utility at every rate, which bounds saving above and makes the cap slack under a single inequality
$2\beta\bar y<(1-\beta R)\bar a$ (Section~\ref{sec:mpc}).

Theorem 2 is a coupling: the two economies are run from a common initial distribution, the higher
rate's policy dominates pointwise, the Markov operator preserves the order, and the order passes to
the limit. The limit exists and is the stationary distribution because of a Doeblin condition that
the borrowing constraint supplies, an atom at zero assets after a run of bad shocks; the same
atom gives uniqueness of the stationary distribution at each rate, which is what turns capital
supply into a function of the rate (Section~\ref{sec:doeblin}).

Theorem 3 is elementary once supply is a monotone function, but Aiyagari's equilibrium condition is
stated at the equilibrium wage, and the household problem is solved at wage one. The bridge is
homogeneity of the stationary distribution in the wage, proved at the level of measures rather than
policies (Section~\ref{sec:homog}), so that the condition $K(r)=\E_r[a]$ at wage $w(r)$ becomes
$k(r)/w(r)=\E^{1}_r[a]$ for the unit-wage economy, and $k/w=\alpha/((1-\alpha)(r+\delta))$ needs no
root.

A last feature deserves mention because it shaped the whole development. Nothing is differentiated.
The value function of a constrained household is not known to be differentiable where the
constraint binds, and the standard envelope theorems do not apply there. The Euler inequalities
that drive every argument are derived from one-sided perturbations of the optimal plan and secant
slopes, in the manner of \citet{ClausenStrub2020}; Lipschitz bounds are secant bounds; monotonicity
is proved by comparing two plans, never by a derivative. This is not a stylistic choice: it is what
allowed the proofs to be completed at all in a setting where a proof assistant will not let a
differentiability assumption pass unnoticed.

\paragraph{Higher risk aversion.}
For $\mu\in\{3,5\}$ the argument fails at Theorem 1, and we show exactly how. The Euler-chain proof
of Theorem 1 reduces to a single \emph{elasticity condition} on the consumption function
(Section~\ref{sec:elasticity}): at tomorrow's assets $b$,
$R_1u'(c_1(b))\le R_2u'(c_2(b))$, that is, for CRRA, $(c_2/c_1)^\mu\le R_2/R_1$. For $\mu>1$ this
is false at high wealth, by the asymptotic linearity of the consumption function
\citep{MaToda2022}: the asymptotic marginal propensity to consume rises with the rate when $\mu>1$,
so consumption of the very rich rises with the rate by more than the condition allows. At
Aiyagari's numbers the condition holds for households with wealth up to about thirteen times mean
income at $\mu=3$ and seven times at $\mu=5$, and fails beyond. Saving of the rich nonetheless
rises with the rate; the sufficient condition simply does not see it.

The chain, it turns out, throws a term away. It starts from the supposition that saving falls with
the rate, which gives $c_2(a)>c_1(a)+(r_2-r_1)a$, and uses only $c_2(a)>c_1(a)$. Keeping the extra
interest income weakens the condition to $(c_2(b)/c_1(b))^\mu\le
(R_2/R_1)(1+(r_2-r_1)a/c_1(a))^\mu$ (Section~\ref{sec:slack}), and in the high-wealth limit this
holds for every $\mu$, exactly, with a margin in logarithmic terms of $1/(\kappa R)$ per unit of
rate difference, about $25$ at Aiyagari's parameters. The slack condition is therefore numerically
true at every wealth level. What we cannot do is prove it from primitives: it asks for the ratio of
consumption at two rates to a relative precision of the order of the rate difference, uniformly in
wealth, and no bound we have on the consumption function is that sharp. That is the frontier, and
we state it as such.

\paragraph{Huggett.}
In \citet{Huggett1993} the asset is a bond in zero net supply and households borrow. Against a
constant supply, monotone capital supply is not enough for single crossing, strict monotonicity is
needed, and we prove the strict form of Theorem 1 for creditors. For debtors a higher rate lowers
cash on hand, and saving can fall with the rate by up to the extra interest on the debt. We carry
this loss through the coupling and obtain $K_1\le K_2+(r_2-r_1)|\underline a|/(1-\Thorn_2)$, where
$\Thorn_2=(\beta R_2)^{1/\sigma}$ is the marginal propensity to save of the rich
(Section~\ref{sec:huggett}). At Huggett's numbers the amplification is about $1500$, and we have no
lower bound on the creditors' gain to set against it. We also show that the propensity-to-save
bound the argument needs cannot be obtained on a bounded asset region from the concavity of the
consumption function and its minimal-MPC floor, although it can on a ray, and we give the
counterexample as a theorem.

\paragraph{What is and is not claimed.}
The result is about Aiyagari's model as he wrote it, with one difference: assets are confined to
$[0,\bar a]$ for some cap $\bar a$, which his computation also imposes, and the theorem carries one
inequality saying the cap is large enough not to bind. Everything else is his: the utility, the
technology, the earnings process up to the two structural conditions. The uniqueness is of the
equilibrium interest rate; the stationary distribution at that rate is unique as well, by the
Doeblin argument. Existence is proved on any interval $[r_{lo},r_{hi}]\subset(-\delta,\lambda)$
whose upper end carries enough capital supply to exceed demand; that supply floor is the one
quantitative fact about the log household this paper does not prove, and Section~\ref{sec:exist}
says why. The bisection theorem takes an equilibrium as given and certifies the algorithm.

\paragraph{Related work.}
Uniqueness in Bewley economies is treated by \citet{Acikgoz2018}, \citet{Light2020},
\citet{Light2022}, and surveyed in \citet{TodaWalsh2024}; multiplicity at high risk aversion is
documented by \citet{YoungWalsh2025}. The consumption-function theory we lean on is
\citet{CarrollKimball1996} for concavity, \citet{MaToda2022} for asymptotic linearity and the
minimal MPC, \citet{CarrollShanker2026} for the buffer-stock bounds, and \citet{LiStachurski2014}
for the contraction in marginal-utility space that gives sufficiency of the Euler equation.
The formalisation uses the fixed-point
and monotone-operator results of \citet{SargentStachurski2024}, several of which we formalised
along the way (Section~\ref{sec:aux}). On formal verification in economics we know of no prior
formalisation of a recursive macroeconomic equilibrium; the closest work is the mean-field-games
literature on uniqueness, \citet{LasryLions2007}, \citet{GraberMatter2024},
\citet{CannarsaCapuani2018}, whose monotonicity conditions we discuss in Section~\ref{sec:aux}
and show do not reach stationary discounted equilibria.

Section~\ref{sec:model} sets out the model as formalised. Section~\ref{sec:main} states the
results for $\mu=1$. Section~\ref{sec:arch} explains the proof. Sections~\ref{sec:beyond} and
\ref{sec:huggett} give the frontier for $\mu>1$ and for Huggett. Section~\ref{sec:aux} collects
auxiliary results. Appendix~\ref{app:index} is the theorem index and Appendix~\ref{app:formal}
describes the formalisation.
